% ÉTUDE DE CAS — analyser une situation, expliquer les risques et proposer un plan.
% Copier dans le dossier du devoir ; remplacer les champs [entre crochets].
\documentclass[11pt,a4paper]{article}

\newcommand{\chargerpreambule}[1]{%
  \IfFileExists{../../templates/#1}{\input{../../templates/#1}}{%
  \IfFileExists{../templates/#1}{\input{../templates/#1}}{%
  \IfFileExists{templates/#1}{\input{templates/#1}}{\input{#1}}}}}
\chargerpreambule{preamble-commun.tex}
\chargerpreambule{preamble-etude-de-cas.tex}
\contextedevoir{ACX401.1 — ECNSA}{Chapitre [numéro] — [intitulé]}

\begin{document}

\fairepagegarde{Étude de cas}{[Titre de l’étude]}%
  {[Organisation, système ou situation étudiée]}%
  {[Prénom NOM — classe]}{\today}

\tableofcontents
\newpage

\section{Résumé de l’étude}

En quelques lignes, donner le sujet, les deux ou trois constats majeurs, leur
importance et la décision ou l’action prioritaire. Cette partie s’écrit souvent
en dernier : elle doit résumer ce que le document démontre.

\filconducteur

\section{Contexte et périmètre}

\subsection{Situation étudiée}

[Présenter l’organisation fictive ou réelle, le besoin et les éléments utiles
pour comprendre le cas. Définir les termes techniques peu courants.]

\subsection{Périmètre et limites}

\begin{tabularx}{\linewidth}{>{\bfseries}p{3.2cm} X}
  \toprule
  Élément & Description \\
  \midrule
  Inclus & [Équipements, services, période, données ou processus examinés.] \\
  Hors périmètre & [Éléments non examinés et raison de leur exclusion.] \\
  Sources & [Documents, captures, entretiens ou commandes utilisés.] \\
  Limites & [Ce qui manque pour conclure avec certitude.] \\
  \bottomrule
\end{tabularx}

\subsection{Méthode suivie}

Décrire les étapes dans l’ordre et justifier brièvement les choix :
\begin{enumerate}[leftmargin=*,itemsep=2pt]
  \item [Quelle information a été recueillie, et où ?]
  \item [Comment a-t-elle été vérifiée ou comparée ?]
  \item [Comment les constats ont-ils été classés et priorisés ?]
\end{enumerate}

\section{Constats appuyés par des preuves}

\subsection{Constat 1 — [Intitulé factuel]}

\textbf{Observation.} [Décrire précisément le fait observé, sans supposer sa cause.]

\textbf{Preuve.} [Citer la capture, la commande, la ligne du document ou le
résultat permettant à un lecteur de vérifier le constat.]

\textbf{Explication.} [Traduire le résultat en langage simple : qu’est-ce qui
se passe, et pourquoi est-ce notable ?]

\textbf{Limite.} [Préciser ce que cette preuve ne permet pas, à elle seule,
d’affirmer. Écrire « aucune identifiée » si c’est pertinent.]

% Réutiliser la structure ci-dessus pour chaque constat important.
\subsection{Constat 2 — [Intitulé factuel]}

\textbf{Observation.} [Fait mesurable ou vérifiable.]

\textbf{Preuve et explication.} [Source, résultat puis interprétation simple.]

\section{Risques et priorités}

Un constat décrit une situation observée ; un risque décrit un scénario
possible et ses conséquences. Éviter de présenter une hypothèse comme un fait.

\small
\tableauECNSA
\begin{tabularx}{\linewidth}{>{\bfseries}p{1.9cm} X X >{\raggedright\arraybackslash}p{2.1cm}}
  \toprule
  Priorité & Scénario possible & Conséquence & Justification \\
  \midrule
  Élevée & [Événement, faiblesse exploitée et actif touché.] & [Impact concret.] & [Preuve ou hypothèse à vérifier.] \\
  Moyenne & [Scénario formulé simplement.] & [Impact concret.] & [Raison du classement.] \\
  \bottomrule
\end{tabularx}
\fintableauECNSA
\normalsize

\section{Plan d’action proposé}

\small
\tableauECNSA
\begin{tabularx}{\linewidth}{>{\bfseries}p{1.9cm} X X >{\raggedright\arraybackslash}p{2.2cm}}
  \toprule
  Échéance & Action & Résultat attendu & Responsable \\
  \midrule
  Immédiate & [Mesure simple, prioritaire et vérifiable.] & [Comment vérifier que c’est fait ?] & [Rôle.] \\
  À planifier & [Mesure nécessitant du temps, des moyens ou une validation.] & [Résultat attendu.] & [Rôle.] \\
  \bottomrule
\end{tabularx}
\fintableauECNSA
\normalsize

\section{Conclusion}

Répondre à la question de départ. Rappeler les constats les plus importants,
la priorité retenue et les limites qui restent à lever. Ne pas introduire de
nouveau résultat dans la conclusion.

\section*{Références et annexes}
\addcontentsline{toc}{section}{Références et annexes}

\begin{itemize}[leftmargin=*]
  \item [Référence : titre, auteur ou organisme, date et adresse si nécessaire.]
  \item [Annexe : nom et utilité de la capture, configuration ou donnée jointe.]
\end{itemize}

\end{document}
